% Introducción condensada (PDF pp. 5-7). Se omiten portada, créditos y el recuento histórico detallado.
\section{Qué son los DBA}\label{sec:intro}

Los Derechos Básicos de Aprendizaje (DBA) son un conjunto de \emph{aprendizajes
estructurantes} que los estudiantes deben alcanzar en cada grado escolar. Esta es la
segunda versión (2016) de los DBA de matemáticas, elaborada por el Ministerio de Educación
Nacional (MEN) con la Universidad de Antioquia a partir de la discusión pública de la
primera versión (2015); sigue abierta a revisión.

\fig{intro_p05_1}{Fase 1 del proceso: presentación de la primera versión de los DBA (2015) y revisión con la comunidad educativa (foros regionales y virtuales, documentos académicos, sistematización de la discusión pública).}
\fig{intro_p05_2}{Fase 3 del proceso: presentación de la segunda versión de los DBA (2016) tras la discusión pública (foros regionales y virtuales, documentos académicos, sistematización).}

Un aprendizaje es la conjunción de conocimientos, habilidades y actitudes; es
\emph{estructurante} porque es una unidad básica sobre la cual se edifica el desarrollo
futuro. Los DBA son coherentes con los Lineamientos Curriculares y los Estándares Básicos
de Competencias (EBC), y ayudan a construir rutas de enseñanza año a año para alcanzar los
EBC de cada grupo de grados.

Los DBA \emph{no} son por sí solos una propuesta curricular: deben articularse con los
enfoques, metodologías y contextos de cada institución (PEI, planes de área y de aula).
Aunque se formulan por grado, el docente puede trasladarlos de un grado a otro según los
procesos de sus estudiantes; son aprendizajes amplios que requieren procesos a lo largo
del año, no una sola actividad.

\subsection*{Estructura de cada DBA}
\begin{description}[nosep]
  \item[Enunciado] el aprendizaje estructurante para el área.
  \item[Evidencias de aprendizaje] indicios clave que muestran si se está alcanzando el
    aprendizaje del enunciado.
  \item[Ejemplo] concreta y complementa las evidencias.
\end{description}

\subsection*{Cómo leerlos}
\begin{itemize}[nosep]
  \item La numeración no define un orden de trabajo en el aula: son los aprendizajes que se
    buscan al finalizar el año.
  \item Una misma experiencia de aula puede aportar a varios DBA a la vez.
  \item Las evidencias hacen observable el aprendizaje; algunas se observan pronto y otras
    exigen un proceso más largo.
  \item Los ejemplos muestran lo que el estudiante debe poder hacer según su edad, y pueden
    contextualizarse según la región, las características étnicas y el entorno.
\end{itemize}

\subsection*{Convenciones de esta versión}
Cada grado está en su propio archivo (\texttt{grados/gradoNN.tex}). Las figuras aparecen
en su lugar, como en el original, con una descripción como pie; las imágenes son archivos
PNG en \texttt{figs/} (no están embebidas en el texto) y se indexan en \texttt{figures.csv}.
Se omitieron las ilustraciones puramente decorativas.
